\chapter{Code and Validation of the Prototypes}
\label{anexo:codigo-prototipos}

This appendix explains in detail the code used in the implementation and validation of the prototypes developed throughout the project.
To this end, it includes the most relevant code fragments, the auxiliary scripts used during development and the outputs obtained in the validation tests performed.\\

The complete source code of the project is also available in the public repository \url{https://github.com/nowelito28/TFG}.

\

\

\section{Prototypes, Implementations and Validations}

The main development of the project was structured through several \textit{LKM} prototypes built progressively and based on the functionality validated in the previous ones.
This strategy made it possible to move from a first contact with programming in \textit{Kernel Space} to the construction of the final prototype, capable of generating an authenticated list of the active system processes from that privileged region.

\

\subsection{\textit{Hello World}}
\label{sec:proto-helloworld}

The first development step consisted of building a basic \textit{LKM} whose objective was to verify how the basic workflow associated with \textit{kernel modules} operates.\\

For this purpose, the external compilation file \textit{Makefile}, based on \textit{kbuild}, was initially designed. It delegates the compilation process to the build system of the \textit{kernel} itself and is reused to compile later prototypes.
This file defines the variable \verb|obj-m|, which tells \textit{kbuild} which object file (\texttt{.o}) must be compiled as a loadable module (\texttt{.ko}).
On the other hand, two main rules are declared to structure the compilation flow. The first one, \verb|all|, invokes the build system of the active \textit{kernel} through the path \verb|/lib/modules/$(shell uname -r)/build|, and specifies the path of the module source code to compile in the current directory with \verb|M=$(shell pwd)|. Therefore, in order to compile, the \textit{Makefile} must be located in the same directory as the module source code and compilation must be executed from that directory.
The second rule, \verb|clean|, removes the intermediate files generated during compilation.
Finally, these rules are declared as \textit{phony targets} to avoid conflicts with files of the same name.
In this way, the module can be compiled externally without needing to modify or recompile the system core.

\begin{lstlisting}[language=make]
obj-m := module_name.o

all: aux.sh
	make -C /lib/modules/$(shell uname -r)/build M=$(shell pwd) modules

clean:
	make -C /lib/modules/$(shell uname -r)/build M=$(shell pwd) clean

.PHONY: all clean
\end{lstlisting}
\captionsetup{type=code,justification=centering,singlelinecheck=false}
\captionof{code}[Generic \textit{Makefile} used to compile \textit{LKMs}.]{Generic \textit{Makefile} used to compile \textit{LKMs}.}
\label{cod:makefile}
\vspace{\baselineskip}

If it is necessary to provide a dependency before compilation, such as an auxiliary file, it can be declared in the \verb|all| rule, as shown in the previous code with the file \texttt{aux.sh}.
In this initial prototype it was not necessary to use this feature, so the dependency shown should only be understood as a general example of how to declare it if needed.\\

Once the compilation tool based on \textit{Makefile} and \textit{kbuild} had been defined, a simple module, \verb|hello_world.c|, was implemented in order to validate the basic life cycle of an \textit{LKM}: loading the module into the \textit{kernel}, executing the module load and unload control functions, registering messages in the \textit{kernel} log and later unloading the module.\\

First, the headers required to compile the module were included and the basic metadata of the component were declared:\\

\begin{lstlisting}[language=C]
#include <linux/module.h>
#include <linux/kernel.h>
#include <linux/init.h>

MODULE_LICENSE("Dual BSD/GPL");
MODULE_AUTHOR("Noel Rodriguez Perez");
MODULE_DESCRIPTION("A simple Hello world LKM!");
MODULE_VERSION("0.1");
\end{lstlisting}
\captionsetup{type=code,justification=centering,singlelinecheck=false}
\captionof{code}[Headers and metadata of the \textit{Hello World} module.]{Headers and metadata of the \textit{Hello World} module.}
\label{cod:hello_headers}
\vspace{\baselineskip}

This fragment includes the headers required to program a \textit{Linux LKM} and, on the other hand, defines several descriptive metadata fields for the component, later visible through the \texttt{modinfo} tool.
Among them, the license is especially important because it conditions the behavior of the module inside the \textit{kernel} and determines whether it is marked as \textit{tainted} when loaded, meaning that it is not considered completely clean for debugging and support purposes.
In this case, a dual \textit{BSD/GPL} license is chosen.
On the one hand, compatibility with \textit{GPL}, used by the \textit{Linux kernel} itself, prevents the module from being marked as \textit{tainted}.
On the other hand, the permissive \textit{BSD} license allows the module to be used in proprietary environments without having to share the source code.\\

Next, the module needs to implement its initialization and exit functions:

\begin{lstlisting}[language=C]
static int __init hello_start(void)
{
    printk(KERN_INFO "Loading hello world module...\n");
    printk(KERN_INFO "Hello world!!\n");
    return 0;
}

static void __exit hello_end(void)
{
    printk(KERN_INFO "Goodbye Mr.\n");
}
\end{lstlisting}
\captionsetup{type=code,justification=centering,singlelinecheck=false}
\captionof{code}[Initialization/load and exit/unload functions of the \textit{Hello World} module.]{Initialization/load and exit/unload functions of the \textit{Hello World} module.}
\label{cod:hello_init_exit}
\vspace{\baselineskip}

The function \verb|hello_start| is executed when the module is loaded into the \textit{kernel} with \verb|insmod <module_name>.ko|. In this prototype, its only purpose is to write messages to the \textit{kernel} log \textit{buffer} using the function \verb|printk|.
These messages can be consulted directly through \verb|dmesg| and, depending on the configuration of the logging system, may also appear in \textit{kernel} log files such as \texttt{/var/log/kern.log}.
Similarly, the function \verb|hello_end| is executed when the module is unloaded with \verb|rmmod <module_name>|.\\
The use of \verb|printk| was verified in the source code of the \textit{kernel} message log, exposed in \texttt{include/linux/printk.h}, where the function itself and the different priority levels that can be associated with each trace are defined, such as \verb|KERN_INFO|, \verb|KERN_DEBUG| or \verb|KERN_ERR|.\\

Finally, both handlers must be explicitly associated with the module control flow:

\begin{lstlisting}[language=C]
module_init(hello_start);
module_exit(hello_end);
\end{lstlisting}
\captionsetup{type=code,justification=centering,singlelinecheck=false}
\captionof{code}[Association of the module entry/load and exit/unload points.]{Association of the module entry/load and exit/unload points.}
\label{cod:hello_module_init}
\vspace{\baselineskip}

These macros indicate to the \textit{kernel} which function must be executed during module loading and which one must be invoked during unloading.
With this, the minimum functional structure of an \textit{LKM} is defined and its correct operation can be validated.\\

With administrator privileges (\textit{root}), the module was compiled, loaded into the \textit{kernel}, the traces registered in the \textit{kernel} log were inspected, the module metadata were consulted and, finally, it was unloaded while also verifying execution of the exit routine:

\begin{lstlisting}[language=bash]
$ make
make -C /lib/modules/6.12.41-current-bcm2711/build M=/home/LKM modules
make[1]: Entering directory '/usr/src/linux-headers-6.12.41-current-bcm2711'
  CC [M]  /home/LKM/hello_world.o
  MODPOST /home/LKM/Module.symvers
  CC [M]  /home/LKM/hello_world.mod.o
  LD [M]  /home/LKM/hello_world.ko
make[1]: Leaving directory '/usr/src/linux-headers-6.12.41-current-bcm2711'
$ insmod hello_world.ko
$ dmesg -T | tail
...
[Fri Sep  5 19:45:58 2025] Loading hello world module...
[Fri Sep  5 19:45:58 2025] Hello world!!
$ modinfo hello_world.ko
filename:       /home/LKM/hello_world.ko
version:        0.1
description:    A simple Hello world LKM!
author:         Noel Rodriguez Perez
license:        Dual BSD/GPL
$ rmmod hello_world
$ dmesg -T | tail
...
[Fri Sep  5 19:46:30 2025] Goodbye Mr.
\end{lstlisting}
\captionsetup{type=code,justification=centering,singlelinecheck=false}
\captionof{code}[Compilation, loading, inspection and unloading of the \textit{Hello World} prototype.]{Compilation, loading, inspection and unloading of the \textit{Hello World} prototype.}
\label{cod:hello_validation}
\vspace{\baselineskip}

Through this workflow, it was verified that the module could be correctly compiled as an external \textit{LKM}, loaded without errors into the \textit{kernel} while executing its initialization routine, with its metadata correctly associated, execute its exit routine when unloaded and properly register messages through \verb|printk| during both phases.\\\\

\

\subsection{Virtual File \texttt{/proc/fdev}}
\label{sec:proto-proc-fdev}

The next step consisted of providing the module with a communication interface between \textit{User Space} and \textit{Kernel Space}, implemented in \verb|proc_fdev.c|.
For this purpose, the \texttt{/proc} virtual file system was used to expose \textit{kernel} information and allow controlled interaction with the \textit{LKM} by creating a custom virtual file in this system, \texttt{/proc/fdev}.\\

The main objective of this phase was to verify that the module could create its own entry in \texttt{/proc}, associate read and write handler functions with it and respond in a controlled way to these operations invoked from \textit{User Space}.\\

First, it was necessary to extend the set of headers included in the module in order to work with \texttt{/proc} and with data transfers between \textit{User Space} and \textit{Kernel Space}:

\begin{lstlisting}[language=C]
#include <linux/proc_fs.h>
#include <asm/uaccess.h>
\end{lstlisting}
\captionsetup{type=code,justification=centering,singlelinecheck=false}
\captionof{code}[Additional headers for the prototype based on \texttt{/proc}.]{Additional headers for the prototype based on \texttt{/proc}.}
\label{cod:proc_headers}
\vspace{\baselineskip}

The header \verb|<linux/proc_fs.h>| is incorporated, required to create and manage entries inside the \texttt{/proc} virtual file system, and, on the other hand, \verb|<asm/uaccess.h>|, which imports functionality such as \verb|copy_to_user| and \verb|copy_from_user|, used to transfer information in a controlled and safe way between \textit{kernel} memory and \textit{User-Space} memory.\\

Next, it was necessary to define a reference to the virtual entry and the operations table associated with that entry in \texttt{/proc}:\\

\begin{lstlisting}[language=C]
static struct proc_dir_entry *ent;

static struct file_operations myops =
{
    .owner = THIS_MODULE,
    .read = myread,
    .write = mywrite,
};

static int simple_init(void)
{
    ent = proc_create("fdev", 0660, NULL, &myops);
    if (!ent)
        return -ENOMEM;
    printk(KERN_INFO "/proc/fdev created\n");
    return 0;
}

static void simple_cleanup(void)
{
    proc_remove(ent);
    printk(KERN_INFO "/proc/fdev removed\n");
}

module_init(simple_init);
module_exit(simple_cleanup);
\end{lstlisting}
\captionsetup{type=code,justification=centering,singlelinecheck=false}
\captionof{code}[Reference to the virtual entry and associated operations table in \texttt{/proc}.]{Reference to the virtual entry and associated operations table in \texttt{/proc}.}
\label{cod:proc_entry_ops}
\vspace{\baselineskip}

In this fragment, the pointer \verb|ent| stores the reference to the entry created in \texttt{/proc} when the module is loaded, later allowing it to be managed and removed when the module is unloaded.
On the other hand, the structure \verb|myops| defines the operations table of the virtual file to be created, associating read and write actions on that file with the handler functions \verb|myread| and \verb|mywrite|, respectively.
In addition, the field \verb|.owner = THIS_MODULE| links these operations to the loaded module itself.
Finally, \verb|simple_init| creates the virtual entry \texttt{/proc/fdev} through \verb|proc_create| during module loading, returning \texttt{-ENOMEM} (\textit{no memory error}) if the entry cannot be created, while \verb|simple_cleanup| removes it with \verb|proc_remove| during unloading.\\

The use of \verb|struct proc_dir_entry|, \verb|proc_create| and \verb|proc_remove| was checked against the source code of the \texttt{/proc} file system of the \textit{kernel}, defined in \texttt{include/linux/proc\_fs.h} and developed in the \texttt{fs/proc} subsystem, while \verb|struct file_operations| is implemented in the general file-operations interface exposed in \texttt{include/linux/fs.h}.\\

One aspect worth highlighting in the development of \textit{kernel module} prototypes is the convenience of using \textit{static} to declare variables and functions that are not going to be used outside the module itself, restricting their visibility to the module source code and favoring better code encapsulation. This criterion is maintained throughout the rest of the project.\\

Once the virtual file entry \texttt{/proc/fdev} had been created and its operations table associated, the write and read handlers of that file were implemented:

\begin{lstlisting}[language=C]
enum { BUFSIZE = 100 };

static int irq = 20;
static int mode = 1;

static ssize_t mywrite(struct file *file, const char __user *ubuf,
		       size_t count, loff_t *ppos)
{
    int num, c, i, m;
    char buf[BUFSIZE];

    if (*ppos > 0 || count >= BUFSIZE)
        return -EFAULT;

	printk(KERN_DEBUG "/proc/fdev: write handler\n");

	if (copy_from_user(buf, ubuf, count))
		return -EFAULT;

    num = sscanf(buf, "%d %d", &i, &m);
    if (num != 2)
        return -EFAULT;


    irq = i;
    mode = m;

    c = strlen(buf);
    *ppos = c;

    printk(KERN_DEBUG "/proc/fdev: %d bytes received from userspace\n", c);
    return c;
}

static ssize_t myread(struct file *file, char __user *ubuf,
		      size_t count, loff_t *ppos)
{
    char buf[BUFSIZE];
    int len = 0;

    if (*ppos > 0 || count < BUFSIZE)
        return 0;

    printk(KERN_DEBUG "/proc/fdev: read handler\n");

    len += sprintf(buf, "irq = %d\n", irq);
    len += sprintf(buf + len, "mode = %d\n", mode);

    if (copy_to_user(ubuf, buf, len))
        return -EFAULT;

    *ppos = len;

    printk(KERN_DEBUG "/proc/fdev: %d bytes sent to userspace\n", len);
    return len;
}
\end{lstlisting}
\captionsetup{type=code,justification=centering,singlelinecheck=false}
\captionof{code}[\textit{Handlers} with data transfer between \textit{User Space} and \textit{Kernel Space}.]{\textit{Handlers} with data transfer between \textit{User Space} and \textit{Kernel Space}.}
\label{cod:proc_rw_handlers}
\vspace{\baselineskip}

The write handler, \verb|mywrite|, receives data from \textit{User Space} through \verb|copy_from_user|, moving \verb|count| bytes from the user-memory \textit{buffer} (\verb|ubuf|) to a \textit{buffer} in \textit{kernel} memory (\verb|buf|).
Similarly, \verb|myread| returns data to \textit{User Space} through \verb|copy_to_user|, moving \verb|len| bytes from the \textit{buffer} in \textit{kernel} memory (\verb|buf|) to the user-memory \textit{buffer} (\verb|ubuf|).
In addition, a \textit{single-shot} access semantics was imposed through the file-position pointer (\verb|*ppos|).\\

The use of \verb|copy_from_user| and \verb|copy_to_user| was checked in the \texttt{uaccess} infrastructure of the \textit{kernel} source code, specifically in \texttt{include/linux/uaccess.h}, where both functions rely on the internal primitives \verb|raw_copy_from_user| and \verb|raw_copy_to_user|.\\

In this example, the values \verb|irq| and \verb|mode| are used as integer test variables to validate the exchange of information between \textit{User Space} and \textit{Kernel Space}. This makes it possible to check their correct interpretation from the data received through \verb|sscanf|, as well as their output formatting through \verb|sprintf|, returning errors such as \texttt{-EFAULT} (\textit{bad address}) when the restrictions are not satisfied or when the input format is not the expected one.
On the other hand, an enumerated constant \verb|BUFSIZE| is defined to simply set the maximum size of the temporary \textit{buffer} used, in this case 100 bytes, both for reading and writing.\\

Once the virtual file \texttt{/proc/fdev} and its read and write handlers had been created, the correct operation of the prototype was checked, assuming administrator permissions (\texttt{root}) on the system:

\begin{lstlisting}[language=bash]
$ make
make -C /lib/modules/6.12.41-current-bcm2711/build M=/home/LKM modules
make[1]: Entering directory '/usr/src/linux-headers-6.12.41-current-bcm2711'
  CC [M]  /home/LKM/proc_fdev.o
  MODPOST /home/LKM/Module.symvers
  CC [M]  /home/LKM/proc_fdev.mod.o
  LD [M]  /home/LKM/proc_fdev.ko
make[1]: Leaving directory '/usr/src/linux-headers-6.12.41-current-bcm2711'
$ insmod proc_fdev.ko
$ dmesg -T | tail
...
[Tue Oct 14 17:25:58 2025] /proc/fdev created
$ ls -l /proc/fdev
-rw-rw---- 1 root root 0 Oct  14 17:25 /proc/fdev
$ echo "7 3" > /proc/fdev
$ dmesg -T | tail
...
[Tue Oct 14 17:26:10 2025] /proc/fdev: write handler
[Tue Oct 14 17:26:10 2025] /proc/fdev: 4 bytes received from userspace
$ cat /proc/fdev
irq = 7
mode = 3
$ dmesg -T | tail
...
[Tue Oct 14 17:26:18 2025] /proc/fdev: read handler
[Tue Oct 14 17:26:18 2025] /proc/fdev: 17 bytes sent to userspace

$ rmmod proc_fdev
$ dmesg -T | tail
...
[Tue Oct 14 17:26:30 2025] /proc/fdev removed
\end{lstlisting}
\captionsetup{type=code,justification=centering,singlelinecheck=false}
\captionof{code}[Basic validation of the prototype based on \texttt{/proc/fdev}.]{Basic validation of the prototype based on \texttt{/proc/fdev}.}
\label{cod:proc_validation}
\vspace{\baselineskip}

This validation confirmed that all module functionality executed correctly, both the creation of the virtual entry \texttt{/proc/fdev} and its custom read and write operations, ensuring a basic communication path between \textit{User Space} and \textit{Kernel Space}.\\\\

\

\subsection{\textit{File Handoff}}
\label{sec:proto-file-handoff}

Once a valid entry in \texttt{/proc} had been created and the basic exchange of information between \textit{User Space} and \textit{Kernel Space} had been achieved in a controlled way, the next step consisted of implementing a mechanism that would allow the module to operate on a real file in the file system, materialized in \verb|fd_handoff.c|.
To do this, the transfer of a file reference (\textit{File Descriptor}) from \textit{User Space} to the \textit{LKM} was proposed, using as input interface the same virtual file \texttt{/proc/fdev} developed in the previous prototype.\\

The general objective of this prototype was to verify that the module could receive from \textit{User Space} the \textit{FD} of a file encoded as a string of \textit{bytes} and convert it to an \textit{int} value, internally resolve that reference to locate the corresponding file in the \textit{kernel} and use it to transfer information from that privileged region to a system file.
With this, the administrator can decide in a controlled way which file receives the requested information.\\

It was necessary to extend the module headers again for this development:

\begin{lstlisting}[language=C]
#include <linux/file.h>
#include <linux/fs.h>
#include <linux/kstrtox.h>
#include <linux/security.h>
\end{lstlisting}
\captionsetup{type=code,justification=centering,singlelinecheck=false}
\captionof{code}[Additional headers for the \textit{File Handoff} mechanism.]{Additional headers for the \textit{File Handoff} mechanism.}
\label{cod:file_handoff_headers}
\vspace{\baselineskip}

The header \verb|<linux/file.h>| allows working with internal \textit{kernel} file references, such as the data structure \verb|struct file|, as well as resolving an \textit{FD} through \verb|fget|.
On the other hand, \verb|<linux/fs.h>| provides general definitions from the file subsystem, including opening modes and \textit{flags}.
The header \verb|<linux/kstrtox.h>| is also added, providing safe conversion functions, such as \verb|kstrtoint|, to transform a string of \textit{bytes} into an integer.
Finally, \verb|<linux/security.h>| allows permissions on referenced files to be validated through \textit{kernel} access-control mechanisms, such as \verb|file_permission|.\\

Regarding the read handler, it was simplified with respect to the previous prototype, limiting itself to returning a deterministic information string indicating to the administrator that the module is ready to receive a \textit{File Descriptor} (\textit{FD}) and perform the \textit{File Handoff}:\\\\

\begin{lstlisting}[language=C]
static ssize_t myread(struct file *file, char __user *ubuf,
		      size_t count, loff_t *ppos)
{
    char buf[] = "Waiting for a FD to perform file handoff\n";
    int len = strlen(buf);

    if (*ppos > 0 || count < len)
        return 0;
    printk(KERN_DEBUG "/proc/fdev: read handler\n");

    if (copy_to_user(ubuf, buf, len))
        return -EFAULT;

    *ppos = len;
    printk(KERN_DEBUG "/proc/fdev: %d bytes sent to userspace\n", len);
    return len;
}
\end{lstlisting}
\captionsetup{type=code,justification=centering,singlelinecheck=false}
\captionof{code}[Read function for the \textit{File Handoff} mechanism.]{Read function for the \textit{File Handoff} mechanism.}
\label{cod:file_handoff_read}
\vspace{\baselineskip}

However, the main functionality of this prototype is concentrated in the write handler, which requires auxiliary functions to organize the write logic in a clearer and more modular way:

\begin{lstlisting}[language=C]
static int write_full(struct file *f, const char *buf, int len)
{
	int w = 0;
	loff_t *ppos = &f->f_pos;

	w = kernel_write(f, buf, len, ppos);
	if (w < 0)
		return w;

	if (w != len)
		return -EIO;

	return w;
}
\end{lstlisting}
\captionsetup{type=code,justification=centering,singlelinecheck=false}
\captionof{code}[Auxiliary function for writing to a file from the \textit{LKM}.]{Auxiliary function for writing to a file from the \textit{LKM}.}
\label{cod:file_handoff_write_full}
\vspace{\baselineskip}

The auxiliary function \verb|write_full| encapsulates writing to an already open file from \textit{Kernel Space} through a single call to \verb|kernel_write|.
For this purpose, it uses the pointer associated with the file (\verb|f_pos|) as the write position, checks whether the call returns an error and verifies that the number of written \textit{bytes} (\verb|w|) matches the requested length (\verb|len|).
Otherwise, an error code (\texttt{-EIO}) is returned, keeping the write check centralized inside the \textit{LKM}.\\

The use of \verb|kernel_write| was supported by verifying its implementation inside the \textit{kernel VFS} source code, specifically in the file \texttt{fs/read\_write.c}.\\

\begin{lstlisting}[language=C]
static int val_metadata(struct file *f)
{
    int rv = 0;

    if (!(f->f_mode & FMODE_WRITE) || !(f->f_flags & O_APPEND))
        return -EBADF;

    rv = file_permission(f, MAY_WRITE | MAY_APPEND);
    if (rv < 0)
        return rv;

    return 0;
}
\end{lstlisting}
\captionsetup{type=code,justification=centering,singlelinecheck=false}
\captionof{code}[Auxiliary function for validating file permissions from the \textit{LKM}.]{Auxiliary function for validating file permissions from the \textit{LKM}.}
\label{cod:file_handoff_val_metadata}
\vspace{\baselineskip}

In this case, the auxiliary function \verb|val_metadata| encapsulates the prior validation logic for the file referenced from \textit{User Space}, checking that it can be used safely in \textit{append} mode from \textit{Kernel Space}.
First, it verifies that it was opened with write capability (\verb|FMODE_WRITE|) and with the \verb|O_APPEND| flag, avoiding accesses incompatible with the policy defined for this prototype.
Then, its effective permissions are checked through \verb|file_permission|, ensuring that the \textit{kernel} can perform write and append operations on it through the permissions \verb|MAY_WRITE| and \verb|MAY_APPEND|.
If any of these restrictions is not satisfied, the function returns the corresponding error code, returning \texttt{-EBADF} (\textit{bad file error}) when the opening metadata are not valid.\\

In this part, the \textit{kernel} source code also had to be consulted. The use of \verb|file_permission| was supported by its definition in \texttt{include/linux/fs.h}, while the relationship between \verb|O_APPEND| and \verb|MAY_APPEND| was checked in \texttt{fs/open.c}.
In addition, in \texttt{fs/namei.c}, inside the function \verb|inode_permission|, it is explicitly stated that every check with \verb|MAY_APPEND| must also include \verb|MAY_WRITE|.\\

\begin{lstlisting}[language=C]
static ssize_t mywrite(struct file *file, const char __user *ubuf,
		       size_t count, loff_t *ppos)
{
    int fd = 0;
    int rv = 0;
    char buf[BUFSIZE];
    static const char msg[] = "Kernel-Space: Writing test for file handoff\n";
    struct file *f = NULL;

    if (*ppos > 0 || count > BUFSIZE)
        return -EFAULT;

    printk(KERN_DEBUG "/proc/fdev: write handler\n");

    if (copy_from_user(buf, ubuf, count))
        return -EFAULT;

    rv = strlen(buf);
    *ppos = rv;

    if (kstrtoint(buf, 10, &fd))
        return -EINVAL;

    if (fd < 0)
        return -EBADF;

    f = fget(fd);
    if (!f)
        return -EBADF;

    rv = val_metadata(f);
    if (rv < 0)
        goto out_put;

    rv = write_full(f, msg, strlen(msg));
    if (rv < 0)
        goto out_put;

    printk(KERN_DEBUG "/proc/fdev: %d bytes written to handoff file\n", rv);

out_put:
    fput(f);
    return rv;
}
\end{lstlisting}
\captionsetup{type=code,justification=centering,singlelinecheck=false}
\captionof{code}[Write function for the \textit{File Handoff} mechanism.]{Write function for the \textit{File Handoff} mechanism.}
\label{cod:file_handoff_write}
\vspace{\baselineskip}

Finally, the handler function \verb|mywrite| concentrates the main logic of the \textit{File Handoff} mechanism, receiving from \textit{User Space} the file descriptor encoded as a string of \textit{bytes} through \texttt{/proc/fdev}.
First, as in the previous prototype, it checks that the \textit{single-shot} access semantics is satisfied and that the maximum size set by \verb|BUFSIZE| is not exceeded, then transferring the information to \textit{kernel} memory through \verb|copy_from_user|.\\

Next, the received string of \textit{bytes} is converted to a decimal integer value with \verb|kstrtoint|, obtaining the corresponding \textit{FD}.
In this way, using \verb|fget| and the obtained \textit{FD}, the reference to the real file inside the \textit{kernel} is recovered; its permissions are then verified with \verb|val_metadata| to check that it complies with the established policy, and a static test string is written to it through \verb|write_full|.\\

Finally, the file reference is released with \verb|fput|, returning the corresponding error code in each case if any check fails, with the help of the \verb|goto| label (\verb|out_put|).\\

The use of \verb|kstrtoint|, \verb|fget| and \verb|fput| was checked against the \textit{kernel} source code, locating \verb|kstrtoint| in the file \texttt{lib/kstrtox.c}, while \verb|fget| and \verb|fput| are found in \texttt{fs/file.c} and \texttt{fs/file\_table.c}, with the exact lines varying depending on the \textit{kernel} version considered.\\

To complete this prototype, it was also necessary to implement an additional utility in \textit{User Space}, \texttt{mk\_fhandoff.c}:

\begin{lstlisting}[language=C]
#include <stdlib.h>
#include <stdio.h>
#include <string.h>
#include <err.h>
#include <errno.h>
#include <fcntl.h>
#include <unistd.h>

static int write_full(int fd, const char *buf, int len)
{
	int w = 0;

	w = write(fd, buf, len);
	if (w < 0)
		return -1;

	if (w != len) {
		errno = EIO;
		return -errno;

	}

	return w;
}

int main(int argc, char *argv[])
{
    if (argc != 3) {
        errx(EXIT_FAILURE, "Usage: %s <common_path> <proc_path>\n"
            "e.g.: %s ./file_handoff /proc/fdev", argv[0], argv[0]);
    }

    const char *f_handoff = argv[1];
    const char *f_proc = argv[2];

    int fd_handoff = open(f_handoff, O_CREAT | O_RDWR | O_APPEND, 0666);

    if (fd_handoff < 0) {
        err(EXIT_FAILURE, "Error opening/creating the file %s",
            f_handoff);
    }



    int fd_proc = open(f_proc, O_RDWR);
    if (fd_proc < 0) {
        err(EXIT_FAILURE, "Error opening %s", f_proc);
    }

    char fd_str[12];
    snprintf(fd_str, sizeof(fd_str), "%d", fd_handoff);

    if (write_full(fd_proc, fd_str, strlen(fd_str)) <= 0) {
        close(fd_proc);
        close(fd_handoff);

        err(EXIT_FAILURE, "Error writing in %s", f_proc);
    }

    printf("File descriptor written in %s:\n%d\n", f_proc, fd_handoff);	

    close(fd_proc);
    close(fd_handoff);

    exit(EXIT_SUCCESS);
}
\end{lstlisting}
\captionsetup{type=code,justification=centering,singlelinecheck=false}
\captionof{code}[\textit{User-Space} program \texttt{mk\_fhandoff.c} to activate \textit{File Handoff}.]{\textit{User-Space} program \texttt{mk\_fhandoff.c} to activate \textit{File Handoff}.}
\label{cod:mk_fhandoff}
\vspace{\baselineskip}

This program acts as an intermediary from \textit{User Space}, creating or opening the real file that will act as the destination from the name indicated as the first argument, obtaining its \textit{FD} and sending it to the module through \texttt{/proc/fdev}, whose path is indicated as the second argument when executing the program, thus completing the \textit{File Handoff} mechanism.
It should be noted that this user utility also implements an auxiliary write function (\verb|write_full|), responsible for encapsulating the call to \verb|write| and checking that the operation does not fail and that the number of written \textit{bytes} matches the requested size.\\

Once the complete basic \textit{File Handoff} mechanism had been finished, its operation was validated by checking that the module correctly received the \textit{FD} of a real file from \textit{User Space}, resolved it inside the \textit{kernel} and wrote to that file the string defined in \verb|mywrite|:

\begin{lstlisting}[language=bash]
$ uname -r
6.12.41-current-bcm2711
$ cd LKM
$ make
make -C /lib/modules/6.12.41-current-bcm2711/build M=/home/LKM modules
make[1]: Entering directory '/usr/src/linux-headers-6.12.41-current-bcm2711'
  CC [M]  /home/LKM/fd_handoff.o
  MODPOST /home/LKM/Module.symvers
  CC [M]  /home/LKM/fd_handoff.mod.o
  LD [M]  /home/LKM/fd_handoff.ko
make[1]: Leaving directory '/usr/src/linux-headers-6.12.41-current-bcm2711'
$ insmod fd_handoff.ko
$ dmesg -T | tail
...
[Wed Oct 29 16:25:15 2025] /proc/fdev created
$ cd ..
$ gcc -c -Wshadow -Wvla -Wall mk_fhandoff.c
$ gcc -o mk_fhandoff mk_fhandoff.o
$ cat /proc/fdev
Waiting for a FD to perform file handoff
$ dmesg -T | tail
...
[Wed Oct 29 16:28:27 2025] /proc/fdev: read handler
[Wed Oct 29 16:28:27 2025] /proc/fdev: 41 bytes sent to userspace
$ ./mk_fhandoff ./file_handoff /proc/fdev
File descriptor written in /proc/fdev:
3
$ dmesg -T | tail
...
[Wed Oct 29 16:30:11 2025] /proc/fdev: write handler
[Wed Oct 29 16:30:11 2025] /proc/fdev: 1 bytes written to handoff file
$ cat ./file_handoff
Kernel-Space: Writing test for file handoff

$ cd LKM
$ rmmod fd_handoff

$ dmesg -T | tail
...
[Wed Oct 29 16:32:20 2025] /proc/fdev removed
\end{lstlisting}
\captionsetup{type=code,justification=centering,singlelinecheck=false}
\captionof{code}[Validation of the \textit{File Handoff} mechanism.]{Validation of the \textit{File Handoff} mechanism.}
\label{cod:file_handoff_validation}
\vspace{\baselineskip}

Through this validation, it was confirmed that the module could correctly receive the \textit{FD} of a real file from \textit{User Space}, resolve its reference in the \textit{kernel} and use it as the effective write destination through the \textit{File Handoff} mechanism.\\

As a consequence, a functional channel was implemented to transfer information generated in \textit{Kernel Space} in a controlled way to the file indicated by \textit{User Space}, serving as the basis for the later incorporation of authenticated critical content in the secure report.\\

\

\subsection{\textit{HMAC-SHA256} Cryptographic Authentication Mechanism}
\label{sec:proto-crypto}

Once the \textit{File Handoff} process had been validated, the next step consisted of incorporating into the module \verb|crypto_fd.c| a cryptographic authentication mechanism for the content generated in \textit{Kernel Space} (\textit{HMAC-SHA256}) and including that \textit{HMAC} tag in the final transfer file, in order to guarantee its integrity and authenticity by relying on the previously validated infrastructure.
In this way, the content can later be verified from \textit{User Space} using the same shared secret key.\\

For this purpose, it was necessary to extend the headers required to compile the module and incorporate the embedded symmetric key with which the \textit{HMAC} is calculated, which must also be available to the administrator during later verification.\\

\begin{lstlisting}[language=C]
#include <crypto/hash.h>
#include <linux/base64.h>
#include <linux/crypto.h>

#include "k_embedded.h"
\end{lstlisting}
\captionsetup{type=code,justification=centering,singlelinecheck=false}
\captionof{code}[Additional headers and embedded key for the \textit{HMAC-SHA256} authentication mechanism.]{Additional headers and embedded key for the \textit{HMAC-SHA256} authentication mechanism.}
\label{cod:crypto_headers}
\vspace{\baselineskip}

At this point, \verb|<crypto/hash.h>| and \verb|<linux/crypto.h>| allow the use of \textit{kernel Crypto API} functionality to request a synchronous \verb|shash| instance, while \verb|<linux/base64.h>| allows the binary result of the \textit{HMAC} to be encoded into a portable textual representation in \textit{Base64} format.
On the other hand, including \verb|"k_embedded.h"| introduces an externally generated binary symmetric secret key embedded in the module as a constant vector of \textit{bytes}, so that it can be used directly during the \textit{HMAC} calculation.\\

Consequently, the module compilation process no longer depends only on the main file, but first requires the header containing the embedded key.
For this purpose, an auxiliary script \verb|gen_k_embedded.sh| was designed. It generates a random 64-byte key, stores it in a binary file \verb|key.bin| and then transforms it into the header \verb|k_embedded.h| as a constant vector of \textit{bytes}, making it importable by the \textit{LKM}.
At the same time, the \textit{Makefile} used to compile the \textit{LKM} declares \verb|k_embedded.h| as a prerequisite of the \verb|all| rule.\\

\begin{lstlisting}[language=sh]
#!/bin/sh
set -e

KEY_LEN=64
openssl rand -out key.bin ${KEY_LEN}

xxd -i -n K key.bin | \
  sed -e 's/^unsigned char/const unsigned char/' \
      -e 's/^unsigned int/const unsigned int/' > k_embedded.h

echo "OK: generated k_embedded.h with key K (${KEY_LEN} bytes)."
\end{lstlisting}
\captionsetup{type=code,justification=centering,singlelinecheck=false}
\captionof{code}[\texttt{gen\_k\_embedded.sh} script to generate the \texttt{k\_embedded.h} header with the embedded key.]{\texttt{gen\_k\_embedded.sh} script to generate the \texttt{k\_embedded.h} header with the embedded key.}
\label{cod:gen_k_embedded}
\vspace{\baselineskip}

It should be noted that, in order to use this script once created, it was necessary to assign it execution permissions externally with \verb|chmod +x gen_k_embedded.sh|, allowing it to be executed afterwards with \verb|./gen_k_embedded.sh| in the same working directory.\\

Regarding the programming of the content-encoding logic to obtain the \textit{HMAC} in the \textit{LKM}, this functionality was implemented to calculate \textit{HMAC-SHA256} from a data block treated as an \textit{array} of \textit{bytes}:

\begin{lstlisting}[language=C]
static int get_hmac_sha256(const u8 *buf, int buf_len, u8 **hmac, int *hmac_len)
{
    int rv = 0;
    struct crypto_shash *tfm;

    tfm = crypto_alloc_shash("hmac(sha256)", 0, 0);

    if (IS_ERR(tfm))
        return PTR_ERR(tfm);

    rv = crypto_shash_setkey(tfm, K, K_len);
    if (rv)
        goto out_free_tfm;

    *hmac_len = crypto_shash_digestsize(tfm);

    *hmac = kmalloc(*hmac_len, GFP_KERNEL);
    if (!*hmac) {
        rv = -ENOMEM;
        goto out_free_tfm;
    }

    SHASH_DESC_ON_STACK(desc, tfm);
    desc->tfm = tfm;

    rv = crypto_shash_digest(desc, buf, buf_len, *hmac);

out_free_tfm:
    crypto_free_shash(tfm);
    return rv;
}
\end{lstlisting}
\captionsetup{type=code,justification=centering,singlelinecheck=false}
\captionof{code}[Calculation of \textit{HMAC-SHA256} for a data \textit{byte} string.]{Calculation of \textit{HMAC-SHA256} for a data \textit{byte} string.}
\label{cod:crypto_get_hmac_sha256}
\vspace{\baselineskip}

This function encapsulates the main logic for calculating the \textit{HMAC-SHA256} cryptographic authenticator over an arbitrary data block represented as an \textit{array} of \textit{bytes} in \textit{kernel} memory.\\

To do this, it requests from the \textit{Crypto API} a synchronous \textit{hash} instance (\textit{shash}) associated with the algorithm \verb|hmac(sha256)| through \verb|crypto_alloc_shash|.
Then, the embedded secret key is associated with it through \verb|crypto_shash_setkey| and the \textit{digest} size is obtained with \verb|crypto_shash_digestsize|.
Afterwards, dynamic memory is allocated with \verb|kmalloc| to store the \textit{HMAC} to be calculated with the corresponding length (\verb|hmac_len|). In addition, the macro \verb|SHASH_DESC_ON_STACK| builds on the execution stack the temporary descriptor required by the cryptographic subsystem (\verb|desc|).
After that, \verb|crypto_shash_digest| calculates the complete \textit{HMAC} in a single operation over the received content.\\

Once the raw \textit{HMAC} of the content had been calculated, an auxiliary function was implemented to convert it to \textit{Base64} format so that it could be included in the transfer file in a readable and portable way:

\begin{lstlisting}[language=C]
static int get_hmac_b64(const u8 *hmac, int hmac_len, u8 **hmac_b64,
			int *hmac_b64len)
{
	int hmac_b64cap = BASE64_CHARS(hmac_len);

	*hmac_b64 = kmalloc(hmac_b64cap, GFP_KERNEL);

	if (!*hmac_b64)
		return -ENOMEM;

	*hmac_b64len = base64_encode(hmac, hmac_len, *hmac_b64);
	if (*hmac_b64len < 0) {
		kfree(*hmac_b64);
		return *hmac_b64len;
	}

	return 0;
}
\end{lstlisting}
\captionsetup{type=code,justification=centering,singlelinecheck=false}
\captionof{code}[Conversion of the binary \textit{HMAC} to \textit{Base64}.]{Conversion of the binary \textit{HMAC} to \textit{Base64}.}
\label{cod:crypto_get_hmac_b64}
\vspace{\baselineskip}

This auxiliary function facilitates serialization of the cryptographic authenticator in the transfer file, transforming the previously calculated binary \textit{HMAC} into a textual \textit{Base64} representation.\\

First, it determines the maximum capacity required for the encoded string through the macro \verb|BASE64_CHARS| and dynamically allocates memory to store it.
Then, the function \verb|base64_encode| performs the effective conversion of that raw \textit{HMAC}, returning the resulting string in \verb|hmac_b64| and its final length in \verb|hmac_b64len|.
In this way, the authenticator can be incorporated into the output file in a readable, portable format that can later be easily verified from \textit{User Space}.\\

The operation and implementation of \verb|crypto_alloc_shash|, \verb|crypto_shash_setkey|, \verb|crypto_shash_digestsize|, \verb|crypto_shash_digest| and the macro \verb|SHASH_DESC_ON_STACK| were checked in the synchronous \textit{hash} infrastructure of the \textit{kernel Crypto API}, exposed in \texttt{include/crypto/hash.h} and developed in the \texttt{crypto} subsystem within the source code of the core itself.
On the other hand, encoding the authenticator in \textit{Base64} relied on \texttt{include/linux/base64.h} and on the corresponding implementation in \texttt{lib/base64.c}.\\

At output-format level, the content is no longer written in isolation to the transfer file, but accompanied by an explicit separation between the data generated by the \textit{kernel} and its \textit{HMAC} authentication tag in \textit{Base64} format:

\begin{lstlisting}[language=C]
static const char sep[] = "\n--KERNEL-PS--\n";
static const int sep_len = sizeof(sep) - 1;
static const char sep_hmac[] = "\n--HMAC--\n";
static const int sep_hmac_len = sizeof(sep_hmac) - 1;

static int write_cont_hmac(struct file *f, const char *cont, int cont_len,
			   	const char *hmac_b64, int hmac_b64len)
{
	int w = 0;
	int total = 0;

	w = write_full(f, sep, sep_len);
	if (w < 0)
		return w;
	total += w;

	w = write_full(f, cont, cont_len);
	if (w < 0)
		return w;
	total += w;

	w = write_full(f, sep_hmac, sep_hmac_len);
	if (w < 0)
		return w;
	total += w;

	w = write_full(f, hmac_b64, hmac_b64len);
	if (w < 0)
		return w;
	total += w;

	return total;
}
\end{lstlisting}
\captionsetup{type=code,justification=centering,singlelinecheck=false}
\captionof{code}[Joint writing of the content and its \textit{HMAC} in \textit{Base64}.]{Joint writing of the content and its \textit{HMAC} in \textit{Base64}.}
\label{cod:crypto_write_cont_hmac}
\vspace{\baselineskip}

This function organizes the final serialization of the authenticated content inside the transfer file.
To do this, it sequentially writes distinguishable textual separators, starting with an initial separator before the block generated by the \textit{kernel} (\verb|sep|), the content to be protected, the separator \verb|sep_hmac| and, finally, the \textit{HMAC} already encoded in \textit{Base64}, reusing the \verb|write_full| function validated in the previous prototype.\\

In this way, the resulting file is structured deterministically, facilitating the separation between the main content to be transferred and the \textit{HMAC} required for later validation from \textit{User Space}.\\

Once the development of the authentication mechanism had been explained, the write handler was modified to integrate all unified components:\\

\begin{lstlisting}[language=C]
static int printh(struct file *f)
{
	int rv = 0;

	static const char cont[] =
		"Kernel-Space: Authentic content protected by HMAC-SHA256\n";
	const int cont_len = sizeof(cont) - 1;

	u8 *hmac = NULL;
	int hmac_len = 0;

	u8 *hmac_b64 = NULL;
	int hmac_b64len = 0;

	rv = get_hmac_sha256(cont, cont_len, &hmac, &hmac_len);
	if (rv < 0)
		goto out;

	rv = get_hmac_b64(hmac, hmac_len, &hmac_b64, &hmac_b64len);
	if (rv < 0)
		goto out_free_hmac;

	rv = write_cont_hmac(f, cont, cont_len, hmac_b64, hmac_b64len);
	if (rv < 0)
		goto out_free_hmacs;

	printk(KERN_INFO "printH: file has been written with HMAC\n");

out_free_hmacs:
	kfree(hmac_b64);


out_free_hmac:
	kfree(hmac);

out:
	return rv;
}
\end{lstlisting}
\captionsetup{type=code,justification=centering,singlelinecheck=false}
\captionof{code}[Joint integration of the authentication mechanism.]{Integration of the authentication mechanism.}
\label{cod:crypto_printh}
\vspace{\baselineskip}

Now, the function \verb|printh| concentrates the main integration logic of the mechanism designed in this prototype.
In this test case, a constant content string \verb|cont| is used; its \textit{HMAC-SHA256} is calculated, encoded in \textit{Base64} and everything is written in a well-structured way to the transfer file.\\

\

\begin{lstlisting}[language=C]
static ssize_t mywrite(struct file *file, const char __user *ubuf,
						size_t count, loff_t *ppos)
{
	int fd = 0;
	int rv = 0;
	char buf[BUFSIZE];

	if (*ppos > 0 || count > BUFSIZE)
		return -EFAULT;
	printk(KERN_DEBUG "/proc/fdev: write handler\n");

	if (copy_from_user(buf, ubuf, count))
		return -EFAULT;

	rv = strlen(buf);
	*ppos = rv;

	if (kstrtoint(buf, 10, &fd))
		return -EINVAL;
	if (fd < 0)
		return -EBADF;

	struct file *f = fget(fd);
	if (!f)
		return -EBADF;

	rv = val_metadata(f);
	if (rv < 0)
		goto out_put;



	rv = printh(f);
	if (rv < 0)
		goto out_put;
	printk(KERN_DEBUG "mywrite: printH OK for fd %d (%d bytes written)\n", fd, rv);

out_put:
	fput(f);

	return rv;
}
\end{lstlisting}
\captionsetup{type=code,justification=centering,singlelinecheck=false}
\captionof{code}[Write handler adapted to the content-authentication mechanism.]{Write handler adapted to the content-authentication mechanism.}
\label{cod:crypto_mywrite}
\vspace{\baselineskip}

For its part, \verb|mywrite| keeps the general logic of the previous prototype for receiving and resolving the \textit{FD}, but replaces the simple static write with a call to \verb|printh|.\\

Once the authentication mechanism had been incorporated into the \textit{File Handoff} flow in the \textit{LKM}, it was necessary to develop an additional utility in \textit{User Space} to verify the validity of the generated authenticated content, intended to be used with administrator permissions.
For this purpose, the program \verb|check_hmac.c| was implemented. It deserializes the content of the transfer file and separates the main content from the \textit{HMAC} tag in \textit{Base64} format, making it possible to recalculate the authenticator with the same shared secret key and compare both results to ensure content integrity.\\

First, the headers and constants required to compile this program were added, including the embedded key shared between the \textit{kernel} and the administrator for calculating the \textit{HMAC}:

\begin{lstlisting}[language=C]
#include <stdlib.h>
#include <stdio.h>
#include <string.h>
#include <err.h>
#include <errno.h>
#include <fcntl.h>
#include <unistd.h>
#include <openssl/hmac.h>
#include <openssl/evp.h>
#include <openssl/crypto.h>

#include "./LKM/k_embedded.h"

enum { LEN = 1024 * 40 };

const char sep[] = "--KERNEL-PS--\n";
const char sep_hmac[] = "--HMAC--\n";
\end{lstlisting}
\captionsetup{type=code,justification=centering,singlelinecheck=false}
\captionof{code}[Headers, constants and symmetric key required for verification.]{Headers, constants and symmetric key required for verification.}
\label{cod:check_hmac}
\vspace{\baselineskip}

First, the functions \verb|read_until_separator|, \verb|read_hmac_line| and \verb|extract_data| are responsible for deserializing the content of the transfer file generated through the \textit{File Handoff} mechanism, separating the main content from the \textit{HMAC} tag stored in \textit{Base64} format by detecting the defined separators:

\begin{lstlisting}[language=C]
int read_until_separator(FILE *f, char **content, int *content_len)
{
	char *line = NULL;
	size_t cap = 0;
	int len = 0;

	int sep_found = 0;
	int sep_hmac_found = 0;

	*content = (char *)malloc(LEN);
	if (!(*content)) {
		warnx("malloc failed for the content\n");

		return 1;
	}

	while ((len = getline(&line, &cap, f)) != -1) {

		if (!sep_found) {
			if (strcmp(line, sep) == 0) {
				sep_found = 1;
			}

			continue;
		}

		if (strcmp(line, sep_hmac) == 0) {
			sep_hmac_found = 1;

			break;
		}

		if (*content_len + len > LEN) {
			free(line);
			warnx("Content read exceeds 4KB buffer before separator\n");

			return 1;
		}

		memcpy(*content + *content_len, line, len);
		*content_len += len;
	}
	free(line);

	if (len == -1) {
		warnx("Lines reading with buffering failed\n");

		return 1;
	}

	if (*content_len > 0 && (*content)[*content_len - 1] == '\n') {
		*content_len -= 1;
	}


	if (sep_hmac_found == 0) {
		warnx("Separator not found\n");

		return 1;
	}

	return 0;
}

int read_hmac_line(FILE *f, char **hmac_b64, int *hmac_b64_len)
{
	size_t cap = 0;

	*hmac_b64_len = getline(hmac_b64, &cap, f);
	if (*hmac_b64_len == -1) {
		warnx("Error reading HMAC\n");

		return 1;
	}

	return 0;
}

int extract_data(FILE *fhandoff, char **content, int *content_len,
		 char **hmac_b64, int *hmac_b64_len)
{
	if (read_until_separator(fhandoff, content, content_len)
	    || read_hmac_line(fhandoff, hmac_b64, hmac_b64_len)) {
		return 1;
	}

	return 0;
}
\end{lstlisting}
\captionsetup{type=code,justification=centering,singlelinecheck=false}
\captionof{code}[Deserialization of the authenticated file content in \textit{User Space}.]{Deserialization of the authenticated file content in \textit{User Space}.}
\label{cod:check_hmac_extract}
\vspace{\baselineskip}

Next, \verb|calculate_hmac|, \verb|encode_base64| and \verb|calc_and_encode_hmac| locally recalculate the authenticator from the extracted content and the same secret key used in the module, which is shared with the administrator running this program, using functionality from the \verb|openssl| library to simplify the cryptographic calculations.\\

Specifically, the \verb|HMAC| function receives the corresponding cryptographic algorithm, the key and the content together with the lengths of both variables, allowing it to directly calculate the \textit{HMAC} corresponding to the content.\\

On the other hand, \verb|EVP_ENCODE_LENGTH| and \verb|EVP_EncodeBlock| are mainly used to encode the previously recalculated \textit{HMAC} into \textit{Base64} format.

\begin{lstlisting}[language=C]
int calculate_hmac(const unsigned char *content, int content_len,
		   unsigned char *hmac_calc, unsigned int *hmac_calc_len)
{
	if (!HMAC(EVP_sha256(), K, K_len, content, content_len, hmac_calc,
		  hmac_calc_len)) {
		warnx("HMAC calculation failed\n");

		return 1;
	}

	return 0;
}

int encode_base64(unsigned char *hmac_calc, unsigned int hmac_calc_len,
		  unsigned char **hmac_b64_calc, int *hmac_b64_calc_len)
{
	*hmac_b64_calc_len = EVP_ENCODE_LENGTH(hmac_calc_len);
	*hmac_b64_calc = (unsigned char *)malloc(*hmac_b64_calc_len + 1);
	if (!*hmac_b64_calc) {
		warnx("malloc failed for HMAC encoded in base 64\n");

		return 1;
	}

	*hmac_b64_calc_len = EVP_EncodeBlock(*hmac_b64_calc, hmac_calc, hmac_calc_len);
	if (*hmac_b64_calc_len <= 0) {
		warnx("EVP_EncodeBlock failed or returned zero length");

		return 1;
	}

	return 0;
}

int calc_and_encode_hmac(const char *content, int content_len,
			 unsigned char **hmac_b64_calc, int *hmac_b64_calc_len)
{
	unsigned char hmac_calc[EVP_MAX_MD_SIZE];
	unsigned int hmac_calc_len = 0;

	if (calculate_hmac((const unsigned char *)content, content_len,
			   hmac_calc, &hmac_calc_len)) {
		return 1;
	}

	if (encode_base64(hmac_calc, hmac_calc_len, hmac_b64_calc,
			  hmac_b64_calc_len)) {
		return 1;
	}

	return 0;
}
\end{lstlisting}
\captionsetup{type=code,justification=centering,singlelinecheck=false}
\captionof{code}[Local \textit{HMAC} calculation and \textit{Base64} encoding.]{Local \textit{HMAC} calculation and \textit{Base64} encoding.}
\label{cod:check_hmac_calc}
\vspace{\baselineskip}

Finally, \verb|verify_hmac| compares the \textit{HMAC} extracted from the file with the one recalculated directly in \textit{Base64} using \verb|CRYPTO_memcmp|.
This whole set of functionalities is integrated into a coordinated workflow in the \verb|main| function itself, from opening the file to the final validation of the authenticated content:

\begin{lstlisting}[language=C]
int verify_hmac(const char *hmac_b64, int hmac_b64_len,
		unsigned char *hmac_b64_calc, int hmac_b64_calc_len)
{
	if (CRYPTO_memcmp(hmac_b64_calc, hmac_b64, hmac_b64_calc_len)) {
		warnx("\nHMAC verification failed: Invalid HMAC\n");

		return 1;
	}

	printf("\nHMAC verification successful: Valid HMAC\n");
	return 0;
}

int main(int argc, char *argv[])
{
	if (argc != 2) {
		errx(EXIT_FAILURE, "Usage: %s <path_file_handoff>\n"
		     "e.g.: %s ./file_handoff", argv[0], argv[0]);
	}

	const char *path_fhandoff = argv[1];

	FILE *fhandoff = fopen(path_fhandoff, "r");
	if (!fhandoff) {
		warnx("fopen failed\n");

		goto err_fhandoff;
	}

	char *content = NULL;
	int content_len = 0;
	char *hmac_b64 = NULL;
	int hmac_b64_len = 0;

	if (extract_data(fhandoff, &content, &content_len,
			 &hmac_b64, &hmac_b64_len)) {
		goto free_cont;
	}

	fclose(fhandoff);
	fhandoff = NULL;

	printf("Extracted kernel content:\n%s\n", content);
	printf("Extracted HMAC(Base64):\n%s\n", hmac_b64);

	unsigned char *hmac_b64_calc = NULL;
	int hmac_b64_calc_len = 0;

	if (calc_and_encode_hmac(content, content_len,
				 &hmac_b64_calc, &hmac_b64_calc_len)) {
		goto free_all;
	}
	free(content);
	content = NULL;

	printf("Calculated HMAC(Base64): \n%s\n", hmac_b64_calc);
	int ext = verify_hmac(hmac_b64, hmac_b64_len, hmac_b64_calc,
			      hmac_b64_calc_len);

	free(hmac_b64);
	free(hmac_b64_calc);
	exit(ext);

free_all:
	if (hmac_b64_calc) {
		free(hmac_b64_calc);
	}

	if (hmac_b64) {
		free(hmac_b64);
	}


free_cont:
	if (content) {
		free(content);
	}

err_fhandoff:
	if (fhandoff) {
		fclose(fhandoff);
	}

	exit(EXIT_FAILURE);
}
\end{lstlisting}
\captionsetup{type=code,justification=centering,singlelinecheck=false}
\captionof{code}[Authenticator comparison and main verifier flow.]{Authenticator comparison and main verifier flow.}
\label{cod:check_hmac_verify_main}
\vspace{\baselineskip}

Once the prototype was completed, its correct operation was validated by checking that the module could correctly generate authenticated content in \textit{Kernel Space}, write it to the file received through the \textit{File Handoff} mechanism and subsequently verify its integrity from \textit{User Space} with the same secret key shared between the administrator and the \textit{kernel}:

\begin{lstlisting}[language=bash]
$ uname -r
6.12.41-current-bcm2711
$ cd LKM
$ chmod +x gen_k_embedded.sh
$ ./gen_k_embedded.sh
OK: generated k_embedded.h with key K (64 bytes).
$ cat k_embedded.h
const unsigned char K[] = {
        0x4a, 0xd3, 0x21, 0xda, 0xc3, 0x6c, 0x38, 0x13, 0xb1, 0x70, 0x48, 0xac,
        0xb2, 0x6b, 0xf7, 0xe4, 0x1b, 0x7d, 0x56, 0xf2, 0xfe, 0x69, 0xb4, 0x59,
        0x0e, 0x31, 0xd6, 0xa7, 0xf1, 0xee, 0x40, 0x8a, 0x17, 0x97, 0x4e, 0xb0,
        0x7a, 0x40, 0x39, 0x40, 0x3f, 0x29, 0x27, 0xd8, 0xa9, 0x84, 0xda, 0x01,
        0x12, 0x70, 0x23, 0xdd, 0x87, 0x02, 0x1b, 0xb9, 0xaf, 0x8c, 0xee, 0xe2,
        0x90, 0x7c, 0x00, 0x10
};

const unsigned int K_len = 64;
$ make
make -C /lib/modules/6.12.41-current-bcm2711/build M=/home/LKM modules
make[1]: Entering directory '/usr/src/linux-headers-6.12.41-current-bcm2711'
  CC [M]  /home/LKM/crypto_fd.o
  MODPOST /home/LKM/Module.symvers
  CC [M]  /home/LKM/crypto_fd.mod.o
  LD [M]  /home/LKM/crypto_fd.ko
make[1]: Leaving directory '/usr/src/linux-headers-6.12.41-current-bcm2711'
$ insmod crypto_fd.ko
$ cat /proc/fdev
Waiting for a FD to perform file handoff
$ dmesg -T | tail
...
[Fri Nov 14 18:25:38 2025] /proc/fdev created
...
[Fri Nov 14 18:28:52 2025] /proc/fdev: read handler
[Fri Nov 14 18:28:52 2025] /proc/fdev: 41 bytes sent to userspace
$ cd ..
$ gcc -c -Wshadow -Wvla -Wall mk_fhandoff.c
$ gcc -o mk_fhandoff mk_fhandoff.o
$ gcc -c -Wshadow -Wvla -Wall check_hmac.c
$ gcc -o check_hmac check_hmac.o -lcrypto
$ ./mk_fhandoff ./file_handoff /proc/fdev
File descriptor written in /proc/fdev:
3
$ dmesg -T | tail
...
[Fri Nov 14 18:31:22 2025] /proc/fdev: write handler
[Fri Nov 14 18:31:22 2025] printH: file has been written with HMAC
[Fri Nov 14 18:31:22 2025] mywrite: printH OK for fd 3 (119 bytes written)
$ cat ./file_handoff

--KERNEL-PS--
Kernel-Space: Authentic content protected by HMAC-SHA256

--HMAC--
xUPsVlgNMcm1Uxj2mac0uTD3gdkDTyeNvbXuBZjq6Bw=

$ ./check_hmac ./file_handoff
Extracted kernel content:
Kernel-Space: Authentic content protected by HMAC-SHA256

Extracted HMAC(Base64):
xUPsVlgNMcm1Uxj2mac0uTD3gdkDTyeNvbXuBZjq6Bw=
Calculated HMAC(Base64):
xUPsVlgNMcm1Uxj2mac0uTD3gdkDTyeNvbXuBZjq6Bw=

HMAC verification successful: Valid HMAC
$ cd LKM
$ rmmod crypto_fd
$ dmesg -T | tail
...
[Fri Nov 14 18:31:47 2025] /proc/fdev removed
\end{lstlisting}
\captionsetup{type=code,justification=centering,singlelinecheck=false}
\captionof{code}[Validation of the authentication mechanism using \textit{HMAC-SHA256}.]{Validation of the authentication mechanism using \textit{HMAC-SHA256}.}
\label{cod:crypto_validation}
\vspace{\baselineskip}

This validation confirmed that the module could correctly generate authenticated content in \textit{Kernel Space}, attach its \textit{HMAC-SHA256} tag in \textit{Base64} format and transfer the complete set to the real file received through the \textit{File Handoff} mechanism.
Likewise, the \verb|check_hmac.c| program made it possible to verify from \textit{User Space} that the authentication tag received in the file matched the one recalculated locally from the extracted content and the same shared secret key, thereby correctly validating the operation of the content-authenticity mechanism designed in this prototype.\\\\

\

\subsection{Secure Reporting System for Running System Processes}
\label{sec:proto-report}

Finally, once the authentication mechanism for the content to be transferred using \textit{HMAC-SHA256} had been validated, the next step consisted of replacing the fixed test string with real content that is genuinely critical for the system and generated dynamically in \textit{Kernel Space}.
For this purpose, a global list of all running system processes was implemented, with a format inspired by the output of \texttt{ps}, so that this content could be transferred to the handoff file in an authenticated way and subsequently verified from \textit{User Space} in the same way as in the previous prototype.
This module was designed in the \verb|crypto_ps.c| file.\\

At this point, all the infrastructure developed in the previous prototypes is reused without substantial changes, including the \textit{File Handoff} mechanism, the joint writing of the content and its \textit{HMAC} in \textit{Base64}, as well as the utilities for writing to the \texttt{/proc} file and verifying the result in \textit{User Space}.\\

Initially, new constants and headers had to be incorporated into the module:

\begin{lstlisting}[language=C]
#include <linux/cred.h>
#include <linux/pid.h>
#include <linux/rcupdate.h>
#include <linux/security.h>
#include <linux/sched.h>
#include <linux/nsproxy.h>
#include <linux/pid_namespace.h>
enum {
	MAX_SIZE = 1024 * 40,
	TIPIC_HMACB64_SIZE = 44,
	UID_SIZE = 12,
	UIDNS_SIZE = 14,
	PID_SIZE = 12,
	PIDNS_SIZE = 14,
	GID_SIZE = 11,
	CMD_SIZE = 50,
	PS_LINE_SIZE = UID_SIZE + UID_SIZE + UIDNS_SIZE +
	    PID_SIZE + PID_SIZE + PIDNS_SIZE + GID_SIZE + CMD_SIZE,
};
static const char header[] =
    "UID_KERNEL  UID_LOCAL   UID_NS        PID_KERNEL  PID_LOCAL   PID_NS        GID        COMMAND\n";
static const int header_len = sizeof(header) - 1;
\end{lstlisting}
\captionsetup{type=code,justification=centering,singlelinecheck=false}
\captionof{code}[Headers and constants for generating the process list.]{Headers and constants for generating the process list.}
\label{cod:report_headers_consts}
\vspace{\baselineskip}

In order to build a list of all processes running on the system and in its different namespaces, it was necessary to extend the module headers again.
Specifically, \verb|<linux/cred.h>| and \verb|<linux/security.h>| allow the module to access the credentials of each process, \verb|<linux/pid.h>|, \verb|<linux/sched.h>| and \verb|<linux/pid_namespace.h>| provide the definitions required to traverse the global task list and obtain their identifiers in the different namespaces, while \verb|<linux/rcupdate.h>| provides the \textit{RCU} synchronization mechanism required to perform this traversal safely.
On the other hand, the auxiliary constants defined bound the maximum size of the generated content, the size reserved for each serialized field and the basic structure of the generated process table with its corresponding metadata.\\

Once the headers and constants required to structure the content to be transferred had been established, two auxiliary functions were incorporated to format and dynamically and safely build the complete process table inside a single output \textit{buffer}:

\begin{lstlisting}[language=C]
static int pad_str_right(char *buf, int curr_len, int buf_len, char pad_char)
{
	int padding = 0;

	if (curr_len >= buf_len) {
		padding = 0;

		curr_len = buf_len;

	} else {
		padding = buf_len - curr_len;
	}
	
	if (padding > 0) {
		memset(buf + curr_len, pad_char, padding);

		curr_len += padding;
	}

	return buf_len;
}

static int safe_chunk(u8 *dst, int *current_len, char *src, int src_len)
{
	int max_len = MAX_SIZE - *current_len;

	if (max_len < src_len)
		return -ENOSPC;

	memcpy(dst + *current_len, src, src_len);

	*current_len += src_len;

	return src_len;
}
\end{lstlisting}
\captionsetup{type=code,justification=centering,singlelinecheck=false}
\captionof{code}[Auxiliary functions for formatting and safely building serialized content.]{Auxiliary functions for formatting and safely building serialized content.}
\label{cod:report_pad_safe_chunk}
\vspace{\baselineskip}

The \verb|pad_str_right| function pads each textual field with fill characters until a fixed width is reached, thus guaranteeing a homogeneous tabular representation of the metadata of each process in the table.\\

For its part, \verb|safe_chunk| sequentially incorporates each fragment of information into the global \textit{buffer} that will contain the final table, previously checking that there is enough available space before performing the copy.\\

Next, the main function responsible for traversing the global system task list and building the complete table in memory was implemented:

\begin{lstlisting}[language=C]
static int get_ps(u8 **cont, int *cont_len)
{
	struct task_struct *task;

	*cont = (u8 *) kmalloc(MAX_SIZE, GFP_KERNEL);
	if (!*cont) {
		printk(KERN_ERR "get_ps: kmalloc failed\n");

		return -ENOMEM;
	}

	if (safe_chunk(*cont, cont_len, (char *)header, header_len) < 0)
		goto out_fail;

	rcu_read_lock();

	for_each_process(task) {
		if (*cont_len >= MAX_SIZE - PS_LINE_SIZE -
		    	sep_len - sep_hmac_len - TIPIC_HMACB64_SIZE) {
			printk(KERN_WARNING "get_ps: Buffer full -> Truncating\n");
			goto out_fail;
		}

		if (ps_data(task, *cont, cont_len)) {
			printk(KERN_ERR "get_ps: Extracting process data failed\n");

			goto out_fail;
		}

	}

	rcu_read_unlock();

	printk(KERN_INFO "get_ps: Generated ps-like output of %d bytes.\n", *cont_len);

	return 0;

out_fail:
	rcu_read_unlock();

	if (*cont)
		kfree(*cont);

	*cont = NULL;
	*cont_len = 0;

	return -ENOSPC;
}
\end{lstlisting}
\captionsetup{type=code,justification=centering,singlelinecheck=false}
\captionof{code}[Construction of the global process list in \textit{kernel} memory.]{Construction of the global process list in \textit{kernel} memory.}
\label{cod:report_get_ps}
\vspace{\baselineskip}

This code acts as the main core for generating the content to be transferred, initially allocating a fixed maximum-size \textit{buffer} where the complete list will be stored and first incorporating the header previously defined in the constants.\\

It then traverses the global list of active system processes using the \verb|for_each_process| macro under \textit{RCU} read protection, thereby guaranteeing safe access in the presence of multiple concurrent accesses to the critical \verb|task_struct| structures while their metadata is extracted.
During this traversal, the code checks that there is still enough available space before incorporating each new entry, delegating the retrieval and serialization of the fields corresponding to each process to \verb|ps_data|.\\

Finally, if the traversal finishes correctly, it returns the complete generated content in \verb|cont|, which will be used in \verb|printh|, and its total length in \verb|cont_len|.
On the other hand, in case of failure, the allocated memory is released and the corresponding error code is returned.

\begin{lstlisting}[language=C]
static int ps_data(struct task_struct *task, u8 *cont, int *cont_len)
{
	char kuid_str[UID_SIZE];
	char luid_str[UID_SIZE];
	char uidns_str[UIDNS_SIZE];
	char kpid_str[PID_SIZE];
	char lpid_str[PID_SIZE];
	char pidns_str[PIDNS_SIZE];
	char gid_str[GID_SIZE];
	char comm_str[CMD_SIZE];

	int kuid_len = get_kuid_str(task_uid(task), kuid_str);
	if (kuid_len <= 0)
		goto out_fail;

	if (safe_chunk(cont, cont_len, kuid_str, UID_SIZE) < 0)
		goto out_fail;

	int luid_len = get_luid_str(task, luid_str);
	if (luid_len <= 0)
		goto out_fail;

	if (safe_chunk(cont, cont_len, luid_str, UID_SIZE) < 0)
		goto out_fail;

	int uidns_len = get_uidns_str(task, uidns_str);
	if (uidns_len <= 0)
		goto out_fail;

	if (safe_chunk(cont, cont_len, uidns_str, UIDNS_SIZE) < 0)
		goto out_fail;

	int kpid_len = get_kpid_str(task_pid_nr(task), kpid_str);
	if (kpid_len <= 0)
		goto out_fail;

	if (safe_chunk(cont, cont_len, kpid_str, PID_SIZE) < 0)
		goto out_fail;

	int lpid_len = get_lpid_str(task, lpid_str);
	if (lpid_len <= 0)
		goto out_fail;

	if (safe_chunk(cont, cont_len, lpid_str, PID_SIZE) < 0)
		goto out_fail;

	int pidns_len = get_pidns_str(task, pidns_str);
	if (pidns_len <= 0)
		goto out_fail;

	if (safe_chunk(cont, cont_len, pidns_str, PIDNS_SIZE) < 0)
		goto out_fail;

	int gid_len = get_gid_str(task, gid_str);
	if (gid_len <= 0)
		goto out_fail;

	if (safe_chunk(cont, cont_len, gid_str, GID_SIZE) < 0)
		goto out_fail;

	int comm_len = get_command_str(task, comm_str);
	if (comm_len <= 0)
		goto out_fail;

	if (safe_chunk(cont, cont_len, comm_str, comm_len) < 0)
		goto out_fail;

	return 0;

out_fail:
	return 1;
}
\end{lstlisting}
\captionsetup{type=code,justification=centering,singlelinecheck=false}
\captionof{code}[Serialization of the metadata of one process.]{Serialization of the metadata of one process.}
\label{cod:report_ps_data}
\vspace{\baselineskip}

The \verb|ps_data| function encapsulates the construction of a single row of the final table corresponding to the task currently traversed in \verb|for_each_process|, sequentially extracting and serializing the main metadata of each specific process.\\

For this purpose, the extraction of each metadata item was encapsulated in independent functions:\\

Next, several auxiliary functions were implemented to extract and format the different user identifiers associated with each process, distinguishing between the global identifier visible from the \textit{kernel}, its relative value inside the process namespace and the identifier of the namespace itself:

\begin{lstlisting}[language=C]
static int get_kuid_str(kuid_t uid_struct, char *uid_str)
{
	int uid_len = 0;

	int uid = from_kuid(&init_user_ns, uid_struct);

	uid_len = snprintf(uid_str, UID_SIZE, "%d", uid);

	pad_str_right(uid_str, uid_len, UID_SIZE, ' ');

	return uid_len;
}

static int get_luid_str(struct task_struct *task, char *uid_str)
{
	int uid_len = 0;

	const struct cred *cred = get_task_cred(task);
	if (!cred) {
		printk(KERN_ERR "get_luid_str: No credentials found for task\n");

		return -ENOSPC;
	}

	int uid = from_kuid(cred->user_ns, cred->uid);

	put_cred(cred);

	uid_len = snprintf(uid_str, UID_SIZE, "%d", uid);

	pad_str_right(uid_str, uid_len, UID_SIZE, ' ');

	return uid_len;
}

static int get_uidns_str(struct task_struct *task, char *uidns_str)
{
	const struct cred *cred = get_task_cred(task);
	if (!cred) {
		printk(KERN_ERR "get_uidns_str: No credentials found for task\n");

		return -ENOSPC;
	}

	unsigned int uid_ns = cred->user_ns->ns.inum;

	put_cred(cred);

	int uidns_len = snprintf(uidns_str, UIDNS_SIZE, "%u", uid_ns);

	pad_str_right(uidns_str, uidns_len, UIDNS_SIZE, ' ');

	return uidns_len;
}
\end{lstlisting}
\captionsetup{type=code,justification=centering,singlelinecheck=false}
\captionof{code}[User identifiers of each process.]{User identifiers of each process.}
\label{cod:report_uid_fields}
\vspace{\baselineskip}

These functions make it possible to obtain and format the user identifiers (\textit{UIDs}) associated with each process.
Specifically, \verb|get_kuid_str| obtains the global UID visible from the \textit{kernel}, \verb|get_luid_str| obtains the UID relative to the process namespace and \verb|get_uidns_str| obtains the user namespace to which those credentials belong.\\

\begin{lstlisting}[language=C]
static int get_kpid_str(int pid, char *pid_str)
{
	int pid_len = snprintf(pid_str, PID_SIZE, "%d", pid);

	pad_str_right(pid_str, pid_len, PID_SIZE, ' ');

	return pid_len;
}

static int get_lpid_str(struct task_struct *task, char *pid_str)
{
	struct pid_namespace *ns = task_active_pid_ns(task);
	if (!ns) {
		printk(KERN_ERR "get_lpid_str: No PID namespace found for task\n");
		return -ENOSPC;
	}

	int pid = task_pid_nr_ns(task, ns);

	int pid_len = snprintf(pid_str, PID_SIZE, "%d", pid);

	pad_str_right(pid_str, pid_len, PID_SIZE, ' ');

	return pid_len;
}

static int get_pidns_str(struct task_struct *task, char *pidns_str)
{
	struct pid_namespace *pid_ns = task_active_pid_ns(task);
	if (!pid_ns) {
		printk(KERN_ERR "get_pidns_str: No PID namespace found for task\n");
		return -ENOSPC;
	}

	int pidns_len = snprintf(pidns_str, PIDNS_SIZE, "%u", pid_ns->ns.inum);

	pad_str_right(pidns_str, pidns_len, PIDNS_SIZE, ' ');

	return pidns_len;
}
\end{lstlisting}
\captionsetup{type=code,justification=centering,singlelinecheck=false}
\captionof{code}[Process identifiers of each task.]{Process identifiers of each task.}
\label{cod:report_pid_fields}
\vspace{\baselineskip}

In this case, these functions make it possible to obtain and format the process identifiers (\textit{PIDs}) of each task.
In particular, \verb|get_kpid_str| obtains the global PID visible from the \textit{kernel}, \verb|get_lpid_str| retrieves the PID relative to the active process namespace and \verb|get_pidns_str| identifies the process namespace to which the analyzed task belongs.\\

\begin{lstlisting}[language=C]
static int get_gid_str(struct task_struct *task, char *gid_str)
{
	int gid_len = 0;

	const struct cred *cred = get_task_cred(task);
	if (!cred)
		return -ENOSPC;

	int gid = from_kgid(&init_user_ns, cred->gid);
	put_cred(cred);



	gid_len = snprintf(gid_str, GID_SIZE, "%d", gid);

	pad_str_right(gid_str, gid_len, GID_SIZE, ' ');

	return gid_len;
}

static int get_command_str(struct task_struct *task, char *comm_str)
{
	int comm_len = 0;

	if (task->mm == NULL) {
		comm_len = snprintf(comm_str, CMD_SIZE, "[%s]\n", task->comm);

		goto commd_finished;
	}

	struct file *exe_file = task->mm->exe_file;

	if (exe_file) {
		char *path = d_path(&exe_file->f_path, comm_str, CMD_SIZE);

		if (!IS_ERR(path)) {
			comm_len = snprintf(comm_str, CMD_SIZE, "%s\n", path);

			goto commd_finished;
		}
	}

	comm_len = snprintf(comm_str, CMD_SIZE, "%s\n", task->comm);

commd_finished:
	if (comm_len >= CMD_SIZE)
		comm_len = CMD_SIZE - 1;

	return comm_len;
}
\end{lstlisting}
\captionsetup{type=code,justification=centering,singlelinecheck=false}
\captionof{code}[Group identifiers and command name associated with each process.]{Group identifiers and command name associated with each process.}
\label{cod:report_gid_command}
\vspace{\baselineskip}

Finally, these functions complete the process information with the group identifier and the command associated with its execution.
Specifically, \verb|get_gid_str| obtains and formats the global group identifier associated with the process, and \verb|get_command_str| retrieves the command corresponding to its execution, distinguishing between user processes with an accessible executable, \textit{kernel} threads and cases where only the name stored in \verb|task->comm| can be used.\\

Once the dynamic generation of metadata for the active system processes had been integrated into the authenticated write flow to the handoff file, the complete behavior of the final prototype was validated by incorporating container functionality to simulate different namespaces in the system:
\begin{lstlisting}[language=bash,basicstyle=\ttfamily\fontsize{5.6pt}{6.2pt}\selectfont,breaklines=false,columns=fullflexible,keepspaces=true]
$ uname -r
6.12.41-current-bcm2711
$ cd LKM
$ chmod +x gen_k_embedded.sh
$ ./gen_k_embedded.sh
OK: generated k_embedded.h with key K (64 bytes).
$ make
make -C /lib/modules/6.12.41-current-bcm2711/build M=/home/LKM modules
make[1]: Entering directory '/usr/src/linux-headers-6.12.41-current-bcm2711'
  CC [M]  /home/LKM/crypto_ps.o
  MODPOST /home/LKM/Module.symvers
  CC [M]  /home/LKM/crypto_ps.mod.o
  LD [M]  /home/LKM/crypto_ps.ko
make[1]: Leaving directory '/usr/src/linux-headers-6.12.41-current-bcm2711'
$ insmod crypto_ps.ko
$ cat /proc/fdev
LKM ready to receive file descriptors from user.
$ dmesg -T | tail
...
[Thu Nov 27 18:42:11 2025] New proc file created: /proc/fdev
...
[Thu Nov 27 18:42:11 2025] /proc/fdev: read handler
[Thu Nov 27 18:42:11 2025] /proc/fdev myread: read 48 bytes by userspace
$ cd ..
$ gcc -c -Wshadow -Wvla -Wall mk_fhandoff.c
$ gcc -o mk_fhandoff mk_fhandoff.o
$ gcc -g -c -Wshadow -Wvla -Wall check_hmac.c
$ gcc -g -o check_hmac check_hmac.o -lcrypto
$ cat /etc/docker/daemon.json
{
        "userns-remap": "default"
}
$ systemctl restart docker
$ docker run -d --name docker1 debian:stable sleep 10000
7fd1c45a07c0...
$ docker run -d --name docker2 debian:stable sleep 10000
77276c1a7026...
$ docker run -d --name docker3 debian:stable sleep 10000
4a4421829e80...
$ docker ps
CONTAINER ID   IMAGE           COMMAND         CREATED              STATUS          NAMES
4a4421829e80   debian:stable   "sleep 10000"   29 seconds ago       Up 28 seconds   docker3
77276c1a7026   debian:stable   "sleep 10000"   35 seconds ago       Up 34 seconds   docker2
7fd1c45a07c0   debian:stable   "sleep 10000"   About a minute ago   Up 58 seconds   docker1
$ ./mk_fhandoff ./file_handoff /proc/fdev
File descriptor written in /proc/fdev:
3
$ dmesg -T | tail
...
[Thu Nov 27 18:43:49 2025] /proc/fdev: write handler
[Thu Nov 27 18:43:49 2025] /proc/fdev write: written 1 bytes from the user
[Thu Nov 27 18:43:49 2025] get_ps: Generated ps-like output of 21251 bytes.
[Thu Nov 27 18:43:49 2025] printH: file has been written with HMAC
[Thu Nov 27 18:43:49 2025] mywrite: printH OK for fd 3 (21320 bytes written)
$ cat ./file_handoff

--KERNEL-PS--
UID_KERNEL  UID_LOCAL   UID_NS        PID_KERNEL  PID_LOCAL   PID_NS        GID        COMMAND
0           0           4026531837    1           1           4026531836    0          /usr/lib/systemd/systemd
0           0           4026531837    2           2           4026531836    0          [kthreadd]
0           0           4026531837    3           3           4026531836    0          [pool_workqueue_]
...
113         113         4026531837    1440        1440        4026531836    119        /usr/sbin/chronyd
113         113         4026531837    1448        1448        4026531836    119        /usr/sbin/chronyd
0           0           4026531837    1489        1489        4026531836    1000       /usr/libexec/gdm-session-worker
1000        1000        4026531837    1496        1496        4026531836    1000       /usr/lib/systemd/systemd
1000        1000        4026531837    1497        1497        4026531836    1000       /usr/lib/systemd/systemd-executor
...
0           0           4026531837    8211        8211        4026531836    0          /usr/bin/dockerd
0           0           4026531837    8212        8212        4026531836    0          [kworker/1:1]
0           0           4026531837    8239        8239        4026531836    0          [kworker/1:2H]
0           0           4026531837    8500        8500        4026531836    0          /usr/bin/containerd-shim-runc-v2
100000      0           4026532674    8520        1           4026532678    100000     /usr/bin/sleep
0           0           4026531837    8532        8532        4026531836    0          [kworker/2:0]
0           0           4026531837    8596        8596        4026531836    0          /usr/bin/containerd-shim-runc-v2
100000      0           4026532754    8617        1           4026532758    100000     /usr/bin/sleep
0           0           4026531837    8667        8667        4026531836    0          [kworker/2:2H]
0           0           4026531837    8678        8678        4026531836    0          /usr/bin/containerd-shim-runc-v2
100000      0           4026532834    8698        1           4026532838    100000     /usr/bin/sleep
0           0           4026531837    8741        8741        4026531836    0          [kworker/u20:2]
0           0           4026531837    8935        8935        4026531836    0          [kworker/0:0]
0           0           4026531837    8947        8947        4026531836    0          [kworker/u16:3]
0           0           4026531837    8989        8989        4026531836    0          [kworker/3:2]
0           0           4026531837    9033        9033        4026531836    0          /home/noel/Desktop/TFG/Project/mk_fhandoff

--HMAC--
ckvLUckcgmt6RFfdGzuvgnAEgt9rLPJ5q9FYJoFKAW0=

$ ./check_hmac ./file_handoff
Extracted kernel content:
UID_KERNEL  UID_LOCAL   UID_NS        PID_KERNEL  PID_LOCAL   PID_NS        GID        COMMAND
0           0           4026531837    1           1           4026531836    0          /usr/lib/systemd/systemd
0           0           4026531837    2           2           4026531836    0          [kthreadd]
0           0           4026531837    3           3           4026531836    0          [pool_workqueue_]
...
113         113         4026531837    1440        1440        4026531836    119        /usr/sbin/chronyd
113         113         4026531837    1448        1448        4026531836    119        /usr/sbin/chronyd
0           0           4026531837    1489        1489        4026531836    1000       /usr/libexec/gdm-session-worker
1000        1000        4026531837    1496        1496        4026531836    1000       /usr/lib/systemd/systemd
1000        1000        4026531837    1497        1497        4026531836    1000       /usr/lib/systemd/systemd-executor
...
0           0           4026531837    8211        8211        4026531836    0          /usr/bin/dockerd
0           0           4026531837    8212        8212        4026531836    0          [kworker/1:1]
0           0           4026531837    8239        8239        4026531836    0          [kworker/1:2H]
0           0           4026531837    8500        8500        4026531836    0          /usr/bin/containerd-shim-runc-v2
100000      0           4026532674    8520        1           4026532678    100000     /usr/bin/sleep
0           0           4026531837    8532        8532        4026531836    0          [kworker/2:0]
0           0           4026531837    8596        8596        4026531836    0          /usr/bin/containerd-shim-runc-v2
100000      0           4026532754    8617        1           4026532758    100000     /usr/bin/sleep
0           0           4026531837    8667        8667        4026531836    0          [kworker/2:2H]
0           0           4026531837    8678        8678        4026531836    0          /usr/bin/containerd-shim-runc-v2
100000      0           4026532834    8698        1           4026532838    100000     /usr/bin/sleep
0           0           4026531837    8741        8741        4026531836    0          [kworker/u20:2]
0           0           4026531837    8935        8935        4026531836    0          [kworker/0:0]
0           0           4026531837    8947        8947        4026531836    0          [kworker/u16:3]
0           0           4026531837    8989        8989        4026531836    0          [kworker/3:2]
0           0           4026531837    9033        9033        4026531836    0          /home/noel/Desktop/TFG/Project/mk_fhandoff

Extracted HMAC(Base64):
ckvLUckcgmt6RFfdGzuvgnAEgt9rLPJ5q9FYJoFKAW0=
Calculated HMAC(Base64):
ckvLUckcgmt6RFfdGzuvgnAEgt9rLPJ5q9FYJoFKAW0=

HMAC verification successful: Valid HMAC
$ cd LKM
$ rmmod crypto_ps
$ dmesg -T | tail
...
[Thu Nov 27 18:46:06 2025] Proc file deleted: /proc/fdev
\end{lstlisting}
\captionsetup{type=code,justification=centering,singlelinecheck=false}
\captionof{code}[Validation of the secure metadata reporting system for active processes authenticated using \textit{HMAC-SHA256}.]{Validation of the secure metadata reporting system for active processes authenticated using \textit{HMAC-SHA256}.}
\label{cod:report_validation}
\vspace{\baselineskip}

It should be noted that the complete results of \verb|cat ./file_handoff| and \verb|./check_hmac ./file_handoff| are not shown due to the large number of active processes in the system, which generates a very long table.
For this reason, the validation only presents a representative excerpt of the transferred and verified content, with the \textit{HMAC} and the reported sizes corresponding to the complete output generated during the real execution.\\

This validation confirmed that the module could correctly generate a global list of the active system processes from \textit{Kernel Space}.
It was also verified that the list included processes running in other \textit{Namespaces}, as observed by running several containers with \textit{Docker}, showing differences with respect to the metadata of the rest of the processes, mainly in the values related to \verb|UID_LOCAL|, \verb|PID_LOCAL|, \verb|UID_NS| and \verb|PID_NS|.\\

On the other hand, as in the previous prototype, the \verb|check_hmac.c| program correctly validates from \textit{User Space} that the \textit{HMAC-SHA256} tag in \textit{Base64} matches the one recalculated over the complete process table, thereby guaranteeing secure and authenticated reporting of critical system data between both regions of the system.\\
